%! author = lili
%! Date = 4/27/26

% todo beispiele für abhängige ereignisse

% TODO Skript Kapitel 4

\chapter{Unabhängigkeit von Ereignissen (und Bernoulli-Experimente)}
\label{ch:unabhangigkeit-von-ereignissen-und-bernoulli-experimente}\index{Bernoulli-Experiment}
\untit{Vorlesung 16.05.2025 (V05 2025)}
\section{Stochastische Unabhängigkeit von Ereignissen}\label{sec:stochastische-unabhangigkeit-von-ereignissen}
Was bedeutet stochastische Unabhängigkeit?

\paragraph{Intuition:} Zwei Ereignisse $\erga$ und $\ergb$ sind unabhängig, wenn die Information des Eintretens des jeweils
einen keinen Einfluss darauf hat, als wie wahrscheinlich wir ein Eintreten auch des jeweils anderen einschätzen (
Bspw.\ mehrmals hintereinander Würfeln).

\paragraph{Formal:} Um $\erga$ und $\ergb$ unabhängig zu nennen, sollten also gelten \txc{(2025)}
\[
    \pe(\ergb|\erga) = \pe(\ergb)
    \text{ und }
    \pe(\erga|\ergb) = \pe(\erga) \txo{\ = \frac{\pe(\erga \cap \ergb)}{\pe(\ergb)}}
\]
\txc{($\pe(\erga|\ergb)$: die bedingte Wahrscheinlichkeit von $\erga$ gegeben $B$.)}
wobei jede dieser beiden Gleichungen äquivalent ist zu.
\[
    \pe(\erga \cap \ergb) = \pe(\erga) \cdot \pe(\ergb)
\]

Also
\[
    \pe(\erga|\ergb) = \pe(\erga)  = \frac{\pe(\erga \cap \ergb)}{\pe(\ergb)} \Leftrightarrow \pe(\erga) \cdot \pe(\ergb) = \pe(\erga \cap \ergb)
\]

\txc{
    $\pe(\erga \cap \ergb)$: \gls{ws}, dass $\erga$ und $\ergb$ gleichzeitig eintreten}

%%%%%%%%%%%%%%%%%%%%%
% Unabhängigkeit %
%%%%%%%%%%%%%%%%%%%%%

\begin{defi}
    Zwei Ereignisse $\erga, $\ergb$ \subseteq \eraum $ heißen unabhängig, wenn $\pe(\erga \cap \ergb) = \pe(\erga) \cdot \pe(\ergb)$.
\end{defi}
\begin{bsp}
TODO
\end{bsp}
\begin{beo} \txc{(2025)} \\
unabhängig $\neq$ disjunkt!
    \noindent Also sind disjunkte Ereignisse unabhängig, wenn $0 = \pe(\emptyset) = \pe(\erga \cap \ergb) = \pe(\erga) \cdot \pe(\ergb)$, also muss entweder $\pe(\erga)=0$ oder $\pe(\ergb)=0$ sein (mindestens eines der Ereignisse \gls{ws} Null haben):
    \[
        0 = \pe(\emptyset) = \pe(\erga \cap \ergb) = \pe(\erga) \cdot \pe(\ergb) \Leftrightarrow \pe(\erga)=0 \text{ oder } \pe(\ergb)=0
    \]
\end{beo}
\komm{disjunkte Ereignisse: $\erga \cap $\ergb$ = \emptyset \Rightarrow$ $\erga$ und $\ergb$ können nicht gleichzeitig eintreten. \\
unabhängige Ereignisse: $\pe(\erga \cap \ergb) = \pe(\erga) \cdot \pe(\ergb)$ \\
}
\begin{beo} \txc{(2025)} \\
Es gilt: $\erga$, $\ergb$ unabhängig $\Leftrightarrow$ $\erga, \ergb^c$ unabhängig $\Leftrightarrow \erga^c, \ergb^c$ unabhängig.

    \paragraph{Beweis.} \txo{z.z. $\pe(\erga \cap \ergb^c) = \pe(\erga) \cdot \pe(\ergb^c)$} \\
    \texttt{Siehe Folie 9}
\end{beo}
\begin{beo}
    $\erga$ kann unabhängig von sich selbst sein. \txc{(2025)}
    \[
        \pe(\erga) = \pe(\erga \cap \erga) = \pe(\erga) \cdot \pe(\erga) = \pe(\erga)^2 \Leftrightarrow \pe(\erga) \in \kl{0,1}
    \]
    \txc{Umgangssprachlicher gesagt: $\erga$ kann unabhängig von sich selbst sein, wenn seine \gls{ws} entweder 0 oder 1 ist.}
\end{beo}
\begin{defi}
    Unabhängigkeit von $n$ Ereignissen. \txc{(2025)}  \\
    Die Ereignisse $\erga_1 , \dots  , \erga_n \subseteq \eraum $ heißen unabhängig, wenn
    \[
        \pe(\erga_{i_1} \cap \dots \cap \erga_{i_k}) = \pe(\erga_{i_1}) \cdots \pe(\erga_{i_k})
    \]
    für alle $k = 2, \dots  , n$ und alle Auswahlen von $k$ verschiedenen
    Indizes $i_1 , \dots  , i_k \in \{1, \dots  , n\}$. \txc{(Also alle Paare, Tripel, etc. prüfen.)}
\end{defi}
\begin{bsp}
    \texttt{siehe Folie 14--15}
\end{bsp}
\begin{bem} \txc{(2025)} \\
Es genügt nicht, paarweise Unabhängigkeit zu prüfen!
    Sachen können paarweise unabhängig, aber nicht insgesamt unabhängig sein.
\end{bem}
\begin{bem}  \txc{(2025)}
    Es genügt nicht, $\pe(\erga_1 \cap \dots  \cap \erga_n ) = \pe(\erga_1 ) \cdots \pe(\erga_n )$ zu prüfen.
    Es kann insgesamt unabhängig sein, aber nicht paarweise etc.\ unabhängig.
\end{bem}


\section{Unabhängige Nacheinanderausführung von Zufallsexperimenten}\label{sec:unabhangige-nacheinanderausfuhrung-von-zufallsexperimenten}
\txc{(2025)} Sollen Zufallsexperimente $(\eraum _1, \pe_1), \dots, (\eraum _n, \pe_n)$
unabhängig nacheinander ausgeführt werden, so modellieren wir dies mit dem W-Raum $(\eraum , \pe)$ mit
\[
    \eraum  = \eraum _1 \times \dots \times \eraum _n
\]
und – im Fall, dass wir abzählbare $\eraum _i$ haben –
\[
    \wfk(\ele) = \wfk_1(\ele_1) \cdots \wfk_n(\ele_n) \forall \ele = (\ele_, \dots, \ele_n) \in \eraum
\]
Sind nämlich $\erga_1 \subseteq \eraum _1 , \dots  , \erga_n \subseteq \eraum _n$ , so sind bei dieser Wahl
\begin{align*}
    \ergb_1 & = \erga_1 \times \eraum _2 \times \dots \eraum _n \\
    \ergb_2 & = \eraum _1 \times \erga_2 \times \eraum _3 \times \dots \times \eraum _n \\
    \vdots & \\
    \ergb_n & = \eraum _1 \times \dots \times \eraum _{n-1} \times \erga_n
\end{align*}
\txc{$\ergb_1$ bedeutet: Im ersten Experiment tritt $\erga_1$ ein, die anderen sind egal. }
unabhängig, denn für $k = 2, \dots  , n$ und $i_1 , \dots  , i_k \in \kl{1, \dots  , n}$ verschieden gilt:
\begin{equation}
    \label{eq:omegaiabzaehlbar}
    \oub{\pe(\ergb_{i_1} \cap \dots \cap \ergb_{i_k})}{ = \erga_1 \times \erga_2 \times \dots \times \erga_n} = \pe_{i_1} \erga_{i_1} \cdot \dots \cdot \pe_{i_k} \erga_{i_k} = \pe(\ergb_{i_1}) \cdot \dots \cdot \pe(\ergb_{i_k})
\end{equation}

\paragraph{Beweis.} \texttt{Selbstständig (siehe Hinweis Folie 20)}

\begin{bsp}
    \texttt{Folie 21}
\end{bsp}

\noindent Sind die $\eraum _i$ nicht alle abzählbar, so definieren wir $\pe$ auf $\eraum  = \eraum _1 \times \dots \times \eraum _n$
\begin{equation}
    \label{eq:omegainichtabzaehlbar}
    \pe(\erga_1 \times \dots \times \erga_n) = \pe_1(\erga_1) \cdot \dots \cdot \pe_n(\erga_n) \quad \forall \erga_i \subseteq \eraum _i \forall i
\end{equation}
motiviert durch Formel~\eqref{eq:omegaiabzaehlbar}.
\begin{bem}  \txc{(2025)} \\
    Formel~\eqref{eq:omegainichtabzaehlbar} kann immer als Definition verendet werden.
    Falls alle $\pe_i$ durch \gls{ws} $\wfk_i$ gegeben sind, folgt aus Formel~\eqref{eq:omegainichtabzaehlbar}:
    \begin{align*}
        \wfk((\ele_1, \dots, \ele_n)) & = \pe(\kl{(\ele_1, \dots, \ele_n)}) \\
        & = \pe(\kl{\ele_1} \times \kl{\ele_2} \times \dots \times \kl{\ele_n}) \\
        &~\overset{(\ref{eq:omegainichtabzaehlbar})}{=} \pe_1(\kl{\ele_1}) \cdot \pe_2(\kl{\ele_2}) \cdot \dots \cdot \pe_n(\kl{\ele_n}) \\
        & = \wfk_1(\ele_1) \cdot \wfk_2 (\ele_2) \cdot \dots \cdot \wfk_n(\ele_n)
    \end{align*}
\end{bem}

%%%%%%%%%%%%%%%%%%%%%
% Binomialverteilung %
%%%%%%%%%%%%%%%%%%%%%

\section{Die Binomialverteilung}\label{sec:die-binomialverteilung}

\begin{defi}
    \textbf{Einfaches \bernoulli }  \txc{(2025)} \\
    \txc{Ein \bernoulli  ist ein Zufallsexperiment mit genau zwei möglichen Ergebnissen: Einem ''Treffer'' oder einer ''Niete''.} \\
    $\eraum _1 = \kl{0,1}; \ele = 1$ bedeutet Erfolg, $\ele = 0$ bedeutet Misserfolg.  \\
    $\wfk_1(1) = p, \wfk_1(0) = 1-p$ und dabei heißt $p \in [0,1]$ die Erfolgswahrscheinlichkeit.
\end{defi}
\begin{defi}
    \textbf{\bernoulli der Länge $n$} \\
    Wiederhole dasselbe einfache \bernoulli  $n$-mal unabhängig.
    \[
        \eraum _n = \kl{0,1}^n, \wfk_n((\ele_1, \dots, \ele_n)) = p^{\oob{^{\sum_{i=1}^n \ele_i}}{=k}} \cdot (1 - p)^{n - \oob{^{\sum_{i=1}^n \ele_i}}{=k}}
    \]
    \txo{\begin{align*}
             \eraum _n & = \underbrace{\eraum _1 \times \dots \times \eraum _n}_{n-mal} \\
             \wfk_n((\ele_1, \dots, \ele_n)) &=  \wfk_1(\ele_1) \cdot \dots \cdot \wfk_1(\ele_n)
             & \forall \ele \in \eraum _n \\
             &=   p^{\sum_{i=1}^n \cdot \ele_i} \cdot (1 - p)^{n - \sum_{i=1}^n \ele_i }
    \end{align*}
    }
\end{defi}
\begin{defi}
    \textbf{Binomialverteilung \index{Binomialverteilung}}  \txc{(2025)} \\
    Suche \gls{ws} für genau $k$ Erfolge in \bernoulli  der Länge $n$. \[
                                                                     \erga_k = \kl{S_n = k} = \kl{(\ele_1 , \dots  , \ele_n ) \in \kl{0, 1}^n : \oub{\sum_{i=1}^n \ele_i}{S_n(\ele)} = k }
    \]
    $\binom{n}{k}$ Möglichkeiten, $k$ Einsen auf die $n$ Stellen zu verteilen.
    \[
        \Rightarrow \glssymbol{binomv}(k) \coloneqq  \pe(\erga_k) = \binom{n}{k} p^k (1-p)^{n-k} \text{ für alle } k \in \kl{0,\dots,n}
    \]
    $\glssymbol{binomv}(k)$ ist eine Wahrscheinlichkeitsfunktion auf $\eraum  = \kl{0,\dots,n}$, denn
    \[
        \sum^n_{k=0} \glssymbol{binomv}(k)  = \sum^n_{k=0} \pe(\erga_k) = \pe(\dot\cup^n_{k=0} \erga_k) = \pe(\kl{0,1}^n)= 1
    \]
    \glssymbol{binomv} heißt \binomv mit den Parametern $n$ und $p$. \\ \\
    Eine \gls{zv} $X$, deren \gls{verteilung} $\pe_X$ die $\glssymbol{binomv}$-Verteilung ist, heißt
    \textbf{binomialverteilt mit den Parametern $n$ und $p$}.
\end{defi}
\begin{bsp}
    Einfaches \bernoulli  \texttt{Siehe Folie 29}
\end{bsp}
\begin{bsp}
    \bernoulli  der Länge n \texttt{Siehe Folie 29}
\end{bsp}
\begin{bsp}
    Binomialverteilung \index{Binomialverteilung}: \texttt{Siehe Folie 29}
\end{bsp}
\begin{prob}
    Wie wahrscheinlich ist, dass in einem \bernoulli  der Länge n der erste Erfolg im n-ten Versuch eintritt?  \txc{(2025)}
    \[
        \erga_n = \kl{(0,\dots,0,1)}, \quad \pe(\erga_n) \ \txo{ = \wfk_n((0,\dots,0,1))} = p(1 - p)^{n-1}
    \]
\end{prob}

%%%%%%%%%%%%%%%%%%%%%
% Binomialverteilung %
%%%%%%%%%%%%%%%%%%%%%

\begin{defi}
    \textbf{geometrische Verteilung mit Erfolgswahrscheinlichkeit $p$}   \txc{(2025)}\\
    \texttt{Siehe Folie 33}
\end{defi}
\begin{defi}
    \textbf{Unabhängigkeit von \gls{zv}}  \txc{(2025)} \\
    \texttt{Siehe Folie 35}
\end{defi}

\newpage
